\section{Evaluación y despliegue de modelos de recompensa}

\subsection{Benchmark Stack}

El curso enfatiza que el reward learning debe ir acompañado de un sistema de evaluación, pues de lo contrario es como ajustar parámetros dentro de una caja negra. Un benchmark stack efectivo contiene al menos tres capas:

\begin{itemize}
    \item \textbf{Métricas automáticas}: log-likelihood, reward model score, KL divergence, comparados con la policy original para ver si hay alguna mejora
    \item \textbf{Evaluación humana profunda}: conversaciones de turnos largos, OpenAI Arena, Anthropic ANAI benchmark, entre otros, para calibrar continuamente el reward model mediante etiquetas de preferencia humana
    \item \textbf{Pruebas de seguridad}: mediante adversarial probing, red team prompts, low-resource attack attempts, se verifica que el reward model se mantenga consistente en los edge cases
\end{itemize}

\begin{importantbox}{El principio del «ciclo cerrado» de evaluación}
Datos → recompensa → política → evaluación → datos. En cada ronda hay que registrar logs, reproducir las métricas, comparar contra el baseline, y usar los casos fallidos para entrenar la siguiente ronda del reward model. De lo contrario, un reward score que parece ir en aumento fácilmente puede ser reward hacking.

\end{importantbox}

\begin{figure}[H]
\centering
\includegraphics[width=0.92\textwidth]{08_cs224r_reward_learning_2025.pdf}
\caption{El circuito de retroalimentación de evaluación del reward learning pipeline (fuente: lecture slide)}
\end{figure}

\subsection{Resumen de la sección}

El sistema de evaluación debe atender simultáneamente las métricas automáticas, las preferencias humanas y los disparadores de seguridad; solo la combinación de los tres puede asegurar que el modelo de reward learning mejore realmente la experiencia y no simplemente «haga trampa».

